\begin{table}[htbp]
\centering
\small
\caption{Ficha de las implementaciones comparadas.}
\label{tab:fichas-implementacion}
\begin{tabular}{p{2.2cm}p{4cm}p{4cm}p{4cm}}
\hline
Modelo & Especificacion & Ajuste de hiperparametros & Divergencias declaradas \\
\hline
OLS & minimos cuadrados ordinarios & no aplica & ninguna \\
GWR & GWR con penalizacion Ridge, kernel bisquare adaptativo & si, dentro del entrenamiento de cada particion: bandwidth por AICc y lambda por validacion espacial anidada; nunca mirando el conjunto evaluado & Ridge sobre la formulacion clasica, necesario por la colinealidad entre variables de distancia \\
GNNWR & Du et al. (2020): una red estima multiplicadores locales de los coeficientes globales & no: configuracion tomada del articulo original y mantenida fija; no se exploro rejilla sobre estos datos & estandarizacion en vez de minmax; parada temprana con 10\% de validacion; recorte de gradiente. Equivalencia con el paquete de los autores: ver observacion 1.2 \\
SANNWR-adaptado & Ni et al. (2022) adaptado: la entrada combina distancia espacial y atributiva con peso fijo 0,5/0,5 & no; alpha tampoco se ajusto & el diseno original integra ambas proximidades de forma aprendida; aqui es una combinacion lineal fija. Por eso el rotulo -adaptado \\
SANNWR-original & Ni et al. (2022) sin adaptar: un SAPDNN (red 2 -> 4 -> 1 compartida entre pares) aprende a fundir la distancia espacial y la atributiva en una sola metrica & no; la capa oculta se fijo de antemano & ninguna estructural: es el diseno publicado. Se corre para medir que aporta el componente que la adaptacion sustituye (02d\_sannwr\_original.py) \\
RF & bosque aleatorio sobre las mismas covariables, sin coordenadas & no; valores por defecto declarados de antemano & ninguna \\
\hline
\end{tabular}
\end{table}